\documentclass[a4paper,12pt]{article}

\usepackage{float}
\usepackage{upgreek}
\usepackage{amsmath}
\usepackage{graphicx}
\usepackage{hyperref}
\usepackage[a4paper, margin=2.5cm]{geometry}
\usepackage[utf8]{inputenc}
\usepackage[T1]{fontenc}
\usepackage{lmodern}   % Latin Modern font with extended characters
\usepackage{fontspec}
\usepackage[numbers]{natbib}

\renewcommand{\figurename}{Slika}

\begin{document}

    \begin{titlepage}
        \begin{center}
            \includegraphics[width=4cm]{slike/fnm} \\[1.5cm]
            {\LARGE \textbf{Vračanje iz vesolja}} \\[1cm]
            {\small Seminarska naloga pri predmetu Sistemsko mišljenje} \\[1cm]

            {\Large Avtor: Taj Šobot} \\[0.5cm]
            {\Large Mentor: doc. Dr. Vladimir Grubelnik} \\[1.5cm]
            \date{18. 6. 2025}
        \end{center}
    \end{titlepage}

    \flushleft

    \pagebreak
    \renewcommand{\contentsname}{Kazalo}
    \tableofcontents{}
    \pagebreak
    \section{Uvod}
    Vračanje iz vesolja oziroma orbite je zanimiv izziv.
    Da se iz vesolja varno vrnemo, mora vesoljsko plovilo vso kinetično in potencialno energijo, ki jo je pridobilo pri vzletu, izgubiti.
    Da porabimo najmanj goriva, ne bomo plovilo v celoti upočasnjevali z motorji.
    Z malim manevrom bomo vesoljsko plovilo poslali v atmosfero.
    Takrat na plovilo začne delovati sila upora, ki nam do površja zmanjša hitrost tako, da z uporabo padal lahko varno pristanemo.
    Plovilo se začne bistveno segrevati, saj se skoraj vsa kinetična in potencialna energija spremeni v toploto.
    V tej seminarski nalogi se s problemom segrevanja plovila ne bomo ukvarjali, temveč bomo samo modelirali dinamiko sistema.
    Bolj natančneje se bomo osredotočali na sile, ki delujejo na plovilo in kako bodo razne začetne konfiguracije plovila imele vpliv na trajektorijo.

    \section{Matematični model}
    \subsection{Opis raziskovalnega problema}
    Modeliral bomo plovilo po kapsuli, ki je astronavte misij Apollo vrnila na Zemljo (običajno imenovan: "komandni modul").
    Naše plovilo se bo na začetku nahajalo v nizki Zemeljski orbiti blizu orbite Mednarodne vesoljske postaje.
    Plovilo se bo potem upočasnilo ($150 \frac{m}{s}$), tako da bo tirnica segala v gostejšo atmosfero.
    Zanima nas, kako se bo spreminjal pospešek, ki bo deloval na plovilo in kako bo ta odvisen od vpadnega kota plovila.
    Ker je plovilo namenjeno za prenos astronavtov, je dobro vedeti, kakšen g bojo naši astronavti čutili pri vračanju.
    Pričakujemo, da pri večjih vpadnih kotih (plovilo usmerjeno navzgor – vzgon usmerjen navzgor), bo plovilo dlje časa letelo v višji atmosferi, kjer je zrak bolj redek.
    Tako bomo izgubljali hitrost počasneje in v daljšem časovnem obdobju.
    To nam posledično zmanjša pospešek, ki bo deloval na plovilo in astronavte.
    Na koncu leta bo plovilo na določeni višini odprlo padalo, ki bo dokončno ustavilo plovilo, ter zagotovilo varen pristanek.

    \subsection{Spremenljivke in enačbe}
    Simulacija vključuje določene omejitve, pri tem pa bomo upoštevali naslednje predpostavke:
    \begin{enumerate}
        \item Simulacija poteka v dvodimenzionalnem (2D) kartezičnem koordinatnem sistemu.
        \item Središče Zemlje je postavljeno v izhodišče koordinatnega sistema, to je $(0, 0)$.
        \item Zemlja je modelirana kot popolnoma okrogla in negibna, z vidika severnega pola, pri čemer se vrti v nasprotni smeri urinega kazalca.
        \item Vrtenje Zemlje je upoštevano s kotno hitrostjo zraka $2\pi$ na dan pri izračunu upora.
        \item Plovilo se giblje v ravnini ekvatorja.
        \item Za izračun zračnega upora in vzgona uporabljam kvadratni zakon, pri čemer koeficienta upora in vzgona ostajata konstantna.
        \item Modeliran je le sistem plovilo-Zemlja, brez vpliva Lune ali drugih nebesnih teles.
        \item Plovilo ohranja svojo orientacijo na smer vektorja hitrosti, ne vrti se okoli osi in ne spreminja velikost vpadnega kota na smer vektorja hitrosti.
    \end{enumerate}

    \begin{figure}[h]
        \centering
        \includegraphics[width=0.7\textwidth]{slike/Slika1}
        \caption{Shema sistema planet-plovilo}
        \label{fig:Slika1}
    \end{figure}
    Koordinatni sistem je kartezični in je postavljen v opazovalni sistem ozadja oziroma vesolja, koordinatni sistem je neodvisen od postavitve planeta.
    V simulaciji so obravnavane naslednje spremenljivke: položaj plovila, hitrost plovila, pospešek plovila ter sile, ki delujejo nanj.
    Položaj planeta (Zemlje) je definiran z uporabo krajevnega vektorja, ki je postavljen v izhodišče koordinatnega sistema:
    \[
        \vec{r}_z = (0, 0).
    \]

    Položaj plovila je izražen s krajevnim vektorjem, katerega začetna vrednost je poljubno določena:
    \[
        \vec{r}_p.
    \]

    Hitrost plovila je podana kot odvod krajevnega vektorja po času, pri čemer je začetna vrednost tudi poljubno izbrana:
    \[
        \vec{v}_p = \frac{d\vec{r}_p}{dt}.
    \]

    Pospešek plovila je definiran kot odvod vektorja hitrosti po času:
    \[
        \vec{a}_p = \frac{d\vec{v}_p}{dt}.
    \]

    Pospešek je izračunan z drugim Newtonovim zakonom:
    \[
        \sum \vec{F} = m \vec{a}.
    \]

    \( \sum \vec{F} \) predstavlja vsoto vseh sil, ki delujejo na plovilo, \( m \)pa je masa plovila.

    Vsota vseh sil je sestavljena iz sile gravitacije $\vec{F_g}$, sile upora $\vec{F_U}$ in sile vzgona $\vec{F_V}$.
    Lahko zapišemo:
    \[
        \sum \vec{F} = \vec{F_g} + \vec{F_U} + \vec{F_V} = m \vec{a}.
    \]

    Sila gravitacije je izračunana preko Newtonovega zakona gravitacije, kjer je $G$ gravitacijska konstanta,
    $m_1$ in $m_2$ sta masa planeta in masa plovila, $r$ je oddaljenost plovila od centra planeta in $\hat{r_0}$
    enotski vektor, izračunan kot: $\hat{r_0} = \frac{\vec{r_P} - \vec{r_Z}}{|\vec{r_P} - \vec{r_Z}|}$.
    Ker je v našem primeru $\vec{r_Z} = (0,0)$, lahko zapišemo  $\hat{r_0} = \frac{\vec{r_P}}{|\vec{r_P}|}$
    \[
        \vec{F_g} = -G \frac{m_1 m_2}{r^2} \hat{r_0}.
    \]

    Sila upora se izračuna sledeče:
    \[
        \vec{F_U} = -\frac{1}{2} \rho(h) S C_U |\vec{v}_{\text{rel}}|^2 \hat{v}_{\text{rel}}.
    \]

    Sila vzgona pa se izračuna sledeče:
    \[
        \vec{F_V} = \frac{1}{2} \rho(h) S C_V * \alpha |\vec{v}_{\text{rel}}|^2 \hat{v}_{\text{rel,rot90}}.
    \]

    Iz tega sledi:
    \[
        \frac{|\vec{F_V}|}{|\vec{F_U}|} = \frac{C_U}{C_V}*\alpha
    \]

    \begin{itemize}
        \item \( \rho(h) \) – funkcija gostote zraka z višino od površja Zemlje ($h = |\vec{r_P}| - polmer zemlje $).
        \item \( S = S_0 \cos(\alpha) \) – čelni presek glede na vpadni kot \(-30^\circ < \alpha < 30^\circ \),
        \item \( C_u \) – koeficient upora,
        \item \( C_v \) – koeficient vzgona (na $^\circ $ kota $\alpha$),
        \item \( \vec{v}_{\text{rel}} = \vec{v}_{\text{p}} + \vec{v}_{\text{zraka}} \) – relativna hitrost plovila glede na zrak,
        \item \(  \hat{v}_{\text{rel,rot90}} \) – enotski vektor relativne hitrosti plovila glede na zrak zavrten za $90^\circ$ (določi smer sile vzgona).
    \end{itemize}

    Funkcijo $\rho(h)$ je težko modelirati, zato smo uporabili že obstoječ model imenovan NRLMSISE-00 (več na\cite{fig:vir5}),
    ki nam poda gostoto zraka kot funkcija nadmorske višine.
    Bolj specifično smo model konfigurirali, da nam vrne gostoto zraka na dani višini na ekvatorju ob 12:00 na dan 21. marec.\\
    Graf funkcije $\rho(h)$ \ref{fig:GrafRoH}:
    \begin{figure}[H]
        \centering
        \includegraphics[width=0.7\textwidth]{slike/Gostota}
        \caption{Graf gostote zraka v odvisnosti od vi šine}
        \label{fig:GrafRoH}
    \end{figure}

    Na sledeči sliki imamo diagram sil \( \vec{F_g} + \vec{F_u} + \vec{F_v} \), ki delujejo na naš sistem plovila.
    Kot $\alpha$ je v Sliki\ref{fig:Slika2} pozitiven (navzgor obrnjen).
    \begin{figure}[H]
        \centering
        \includegraphics[width=0.7\textwidth]{slike/slika2NOVA}
        \caption{Shema sistema plovila}
        \label{fig:Slika2}
    \end{figure}
    Kapsula bo "držala" plovilo pod kotom $\alpha$ na vektor relativne hitrosti skozi celoten let.
    Kot $\alpha$ omejimo na: $0^\circ < \alpha < 30^\circ$.

    \subsection{Konstante}
    Konstante uporabljene v simulaciji so:
    \begin{itemize}
        \item gravitacijska konstanta $G = 6.6743\times10^{-11}$,
        \item masa Zemlje $M_1 = 5.972\times10^{24} kg$,
        \item polmer Zemlje $r_z = 6.371\times10^{6} m$,
        \item kotna hitrost Zemlje $\omega = 7.2921159\times10^{-5} 1/s$.
    \end{itemize}
    V simulaciji imam še druge konstante, ki se lahko poljubno nastavijo (ponavadi začetni pogoji):
    \begin{itemize}
        \item masa plovila,
        \item začetna pozicija plovila,
        \item začetna hitrost plovila,
        \item površina plovila,
        \item koeficient upora,
        \item vpadni kot.
    \end{itemize}
    Vse ostale kostnate povezane z atmosfero, vsebuje jih model NRLMSISE-00\cite{fig:vir5} in se z njimi ne bom posredno ukvarjal.

    \section{Simulacija}
    \subsection{Začetni pogoji}
    Začetni pogoji pozicije in hitrosti bojo navedene pri vsaki posamezni simulaciji v poglavju "Rezultati".
    Za parametre plovila pa bojo vedno isti in sledeči.
    \subsubsection{Parametri plovila}
    Parametre plovila sem nastavil na:
    \begin{itemize}
        \item masa plovila $m = 5806 kg$,
        \item površina plovila $S_0 = 12.0 m^2$,
        \item koeficient upora $C_U = 1.2$,
        \item koeficient vzgona (na $1^\circ$ kota) $C_V = 0.011$,
        \item površina padala $S_{padala} = 1524.0 m^2$,
        \item koeficient upora padala $C_{padala} = 0.9$.

    \end{itemize}
    Podatki o parametrih plovila so bili pridobljeni iz javno dostopnih virov, vključno z dokumentacijo NASA in spletnimi enciklopedijami, kot je Wikipedija.
    Koeficienta upora in vzgona sem pridobil iz virov:~\cite{fig:vir3},\cite{fig:vir4},\cite{fig:vir6}.
    V teh virih je navedenih več različnih koeficientov, saj se pri nadzvočnih hitrostih te začnejo nezanemarljivo spreminjati.
    Ker je izmerjen koeficient upora bil med $1.0$ in $1.4$, smo uporabili približno povprečno vrednost $1.2$.
    Koeficient vzgona je bil pri $30^\circ$ med $0.30$ in $0.35$, zato sem uporabil približno vrednost $0.33$, torej $0.011$ na $1^\circ$ kota.


    \subsubsection{Izbira vhodnih podatkov}
    Za modeliranje smo izbrali kapsulo Apollo, čeprav ta ni več v uporabi.
    Razlog za to izbiro je obsežna razpoložljivost zgodovinskih podatkov, dokumentacije in rezultatov testiranj, ki so javno dostopni.
    Kapsula Apollo je primer plovila, zasnovanega za vračanje iz luninih orbit, hkrati pa je bila uporabljena tudi za misije v nizki zemeljski orbiti.
    Zaradi bogatega nabora podatkov o letih je mogoče rezultate simulacije neposredno primerjati z realnimi podatki, kar povečuje zanesljivost modela.
    To omogoča, da bomo z enakim modelom lahko simuliral tudi druge kapsule – z zgolj prilagoditvijo začetnih parametrov plovila – in pri tem zagotovili robustnost in natančnost modela.

    \subsubsection{Numerična integracija in implementacija}
    Za numerični integrator smo uporabil Eulerjevo metodo s časovnim korakom $\Delta t = 0,01 \, \text{s}$.
    Simulacijo smo implementirali v programskem jeziku C++, za vizualizacijo rezultatov in izdelavo grafov pa smo uporabili Python z knjižnico Matplotlib.
    Simulacija se zaključi, ko plovilo doseže tla ali ko je dosežen prednastavljen najdaljši čas izvajanja.
    Model omogoča prilagoditev za simulacijo drugih pogojev, kot so vračanja iz višjih orbit, eliptičnih orbit (npr. z Lune) ali hiperboličnih orbit.

    \subsection{Konec simulacije}
    Ko plovilo doseže nadmorsko višino 3000 m, se aktivira padalo.
    Padalo je modelirano tako, da se površina plovila v 10 sekundah postopno poveča od začetne vrednosti (površina plovila) do končne vrednosti (površina plovila in padala).
    To je doseženo z linearno interpolacijo površine v časovnem intervalu 10 s.
    Simulacija se zaključi, ko plovilo doseže nadmorsko višino 0 m.

    \subsection{Robni pogoji}
    Numerični model vsebuje določene omejitve.
    V simulaciji so vse spremenljivke shranjene v podatkovnem tipu \texttt{double}, ki zagotavlja natančnost približno \(\pm 2,2 \times 10^{-16}\).
    Ta omejitev lahko povzroči kumulativne zaokrožitvene napake, zlasti pri dolgotrajnih simulacijah ali ob hitrih spremembah sil.
    Te napake so bile upoštevane pri izbiri časovnega koraka \(\Delta t = 0,01 \, \text{s}\).
    Za večjo natančnost bi lahko uporabili višjeredne integracijske metode, kot je Runge-Kutta, ali še dodatno zmanjšali časovni korak.
    Začetne vrednosti in konstante, kot so masa plovila in gravitacijska konstanta, so bile izbrane tako, da čim bolj zmanjšajo vpliv zaokrožitvenih napak, vendar lahko v izračunih tira in pospeškov še vedno pride do manjših odstopanj.

    Simulacija vključuje tudi nekatere singularnosti, ki lahko povzročijo deljenje z nič.
    Pri računanju gravitacijske sile je treba pozornost nameniti majhnim razdaljam, saj te povzročajo izjemno velike sile, kar vodi do numeričnih napak in nerealističnih rezultatov.

    Ker je planet v modelu obravnavan kot točka, simulacija v vsakem primeru ni smiselna, če se plovilo nahaja pod površjem planeta (pod nadmorsko višino 0 m).

    \pagebreak
    \section{Rezultati}
    \subsection{Vračanje iz nizkozemeljske orbite}
    Plovilo smo postavili v krožno orbito na višino 410 km in s hitrostjo 7667 $\frac{m}{s}$.
    Ta je podobna orbiti Mednarodne vesoljske postaje.
    Nato smo hitrost zmanjšali za 150 $\frac{m}{s}$, kar je postavilo plovilo v eliptično orbito, katere perigeja (najnižja točka orbite) sega v atmosfero.
    Simulacijo je bila zagnana za vpadne kote $\alpha$ od $0^\circ$ do $30^\circ$ po korakih $5^\circ$.
    Začetni pogoji (vektorji pozicije in hitrosti) so:
    \begin{itemize}
        \item začetna pozicija plovila $\vec{r_{p0}} = [r_z + 410 000m, 0m]$,
        \item začetna hitrost plovila $\vec{v_{p0}} = [0\frac{m}{s}, 7517\frac{m}{s}] $.
    \end{itemize}
    Slika~\ref{fig:graf1}, prikazuje graf tira, po katerem je naše plovilo potovalo (odvisnost x od y).
    V vseh grafih barvne poti označujejo posamezne tire pri različnih kotih.
    Siv krog pa v tem grafu predstavlja zemljino površje.

    \begin{figure}[H]
        \centering
        \includegraphics[width=1.0\textwidth]{slike/Graf1}
        \caption{Graf $x(y)$ pri različnih $\alpha$}
        \label{fig:graf1}
    \end{figure}

    Slika~\ref{fig:graf2}, prikazuje graf višine nad površjem zemlje v odvisnosti od časa.
    \begin{figure}[H]
        \centering
        \includegraphics[width=1.0\textwidth]{slike/Graf2}
        \caption{Graf $h(t)$ pri različnih $\alpha$}
        \label{fig:graf2}
    \end{figure}

    Slika~\ref{fig:graf3}, prikazuje graf magnitude hitrosti (relativno na koordinatni sitem) v odvisnosti od časa.
    \begin{figure}[H]
        \centering
        \includegraphics[width=1.0\textwidth]{slike/Graf3}
        \caption{Graf $|\vec{v_p}|(t)$ pri različnih $\alpha$}
        \label{fig:graf3}
    \end{figure}

    Slika~\ref{fig:graf5}, prikazuje graf magnitude hitrosti (relativno na površino zemlje) v odvisnosti od časa.
    Zdelo se mi je pomembno tudi ponazoriti hitrost plovila relativnega na tla, saj je pri pristanku dobro vedeti svojo hitrost.
    V našem primeru plovilo ne preseže 5 $\frac{m}{s}$.
    Plovilo je namenjeno za pristajanje na vodi, zato je ta vrednost sprejemljiva.
    \begin{figure}[H]
        \centering
        \includegraphics[width=1.0\textwidth]{slike/Graf5}
        \caption{Graf $| \vec{v_{pt}}|(t)$ pri različnih $\alpha$}
        \label{fig:graf5}
    \end{figure}

    Slika~\ref{fig:graf4}, prikazuje graf magnitude pospeška v odvisnosti od časa.
    \begin{figure}[H]
        \centering
        \includegraphics[width=1.0\textwidth]{slike/Graf4}
        \caption{Graf $|\vec{a_p}|(t)$ pri različnih $\alpha$}
        \label{fig:graf4}
    \end{figure}
    Iz grafa pospeška lahko odčitamo, da je plovilo, ki je imelo največji vpadni kot ($30^\circ$), doživelo najmanjši maksimalni pospešek okoli $2g$.
    Plovilo z najmanjšim ($0^\circ$) pa največji, ta je imel maksimalni pospešek $8g$.
    Drugi narastek pospeška je zaradi odpiranja padala tik pred pristankom, kjer pospešek ne naraste do $5g$.
    Med realnim letom vpadni kot ni konstanten, saj imajo prava plovila navigacijske računalnike in ljudi, ki lahko spreminjajo orientacijo plovila med letom.
    Ponavadi ta nadzoruje vpadni kot tako, da kapsula pride v določeno pristajalno območje, hkrati pa ohranja varne temperaturne razmere.
    Pri standardnem vračanju iz nizko zemeljske orbite astronavti doživijo med $2g$ do $4g$.
    Naši rezultati se ujemajo z realnimi pogoji.

    \subsection{Vračanje iz lunine orbite}
    Začetni pogoji:
    \begin{itemize}
        \item začetna pozicija plovila $\vec{r_{p0}} = [r_z + 400 000 000m, 0m]$,
        \item začetna hitrost plovila $\vec{v_{p0}} = [0\frac{m}{s}, 175\frac{m}{s}] $,
    \end{itemize}
    Te smo izbrali tako da smo plovilo nastavili na (približno) višino lunine orbite in spreminjali hitrost, dokler trajektorija ni zadela zemljine atmosfere.

    Slika~\ref{fig:graf8}, prikazuje graf tira zelo eliptične orbite (iz lune).
    Slika~\ref{fig:graf8Z}, pa približan graf.
    \begin{figure}[H]
        \centering
        \includegraphics[width=1.0\textwidth]{slike/LUNA0Graf1}
        \caption{Graf $x(y)$ pri različnih $\alpha$}
        \label{fig:graf8}
    \end{figure}
    \begin{figure}[H]
        \centering
        \includegraphics[width=1.0\textwidth]{slike/LUNA0Graf1Z}
        \caption{Graf $x(y)$ pri različnih $\alpha$}
        \label{fig:graf8Z}
    \end{figure}
    Kot lahko vidimo smo plovilo poslali mimo atmosfere (a še vedno dovolj blizu, da opazimo njen efekt).
    Zmanjšajmo velikost začetne hitrosti za ($-0,5 \frac{m}{s}$):

    Spremenjeni parametri:
    \begin{itemize}
        \item začetna pozicija plovila $\vec{r_{p0}} = [r_z + 400 000 000m, 0m]$,
        \item začetna hitrost plovila $\vec{v_{p0}} = [0\frac{m}{s}, 174,5\frac{m}{s}] $,
    \end{itemize}

    Slika~\ref{fig:graf9}, prikazuje graf tira zelo eliptične orbite (iz lune).
    Slika~\ref{fig:graf9Z}, pa približan graf.
    Slika~\ref{fig:graf9ZZ}, pa še bolj približan graf.

    \begin{figure}[H]
        \centering
        \includegraphics[width=1.0\textwidth]{slike/LUNA05Graf1}
        \caption{Graf $x(y)$ pri različnih $\alpha$}
        \label{fig:graf9}
    \end{figure}
    \begin{figure}[H]
        \centering
        \includegraphics[width=1.0\textwidth]{slike/LUNA05Graf1Z}
        \caption{Graf $x(y)$ pri različnih $\alpha$}
        \label{fig:graf9Z}
    \end{figure}
    \begin{figure}[H]
        \centering
        \includegraphics[width=1.0\textwidth]{slike/LUNA05Graf1ZZ}
        \caption{Graf $x(y)$ pri različnih $\alpha$}
        \label{fig:graf9ZZ}
    \end{figure}
    Kot lahko vidimo so se nekatera plovila "odbila" od atmosfere druga pa je že takoj potegnilo do površja.
    To pomeni, da nekatera plovila vstopajo v atmosfero več kot enkrat, kar v pravi odpravi ne bi bilo zaželeno.
    
    \vspace{1cm}

    Gremo še malo zmanjšati velikost začetne hitrosti ($-0,5 \frac{m}{s}$):

    Spremenjeni parametri:
    \begin{itemize}
        \item začetna pozicija plovila $\vec{r_{p0}} = [r_z + 400 000 000m, 0m]$,
        \item začetna hitrost plovila $\vec{v_{p0}} = [0\frac{m}{s}, 174\frac{m}{s}] $,
    \end{itemize}

    Slika~\ref{fig:graf10}, prikazuje graf tira zelo eliptične orbite (iz lune).
    Slika~\ref{fig:graf10Z}, pa približan graf.
    Slika~\ref{fig:graf10ZZ}, pa še bolj približan graf.

    \begin{figure}[H]
        \centering
        \includegraphics[width=1.0\textwidth]{slike/LUNA1Graf1}
        \caption{Graf $x(y)$ pri različnih $\alpha$}
        \label{fig:graf10}
    \end{figure}
    \begin{figure}[H]
        \centering
        \includegraphics[width=1.0\textwidth]{slike/LUNA1Graf1Z}
        \caption{Graf $x(y)$ pri različnih $\alpha$}
        \label{fig:graf10Z}
    \end{figure}
    \begin{figure}[H]
        \centering
        \includegraphics[width=1.0\textwidth]{slike/LUNA1Graf1ZZ}
        \caption{Graf $x(y)$ pri različnih $\alpha$}
        \label{fig:graf10ZZ}
    \end{figure}

    Zdaj pa lahko vidimo, da vse trajektorije končajo na Zemljinem površju po prvem vstopu v atmosfero.


    \pagebreak
    \section{Zaključek}
    V tej seminarski nalogi sem ustvaril model za simulacijo vračanja vesoljskega plovila, modeliranega po Apollo kapsuli, od različnih Zemeljskih orbit do površja
    Osredotočili smo se na dinamiko sistema, zlasti na sile, ki delujejo na plovilo (gravitacijo, zračni upor in vzgon), ter na vpliv vpadnega kota $\alpha$ na trajektorijo, hitrost in pospešek.
    Rezultati pokažejo, da večji vpadni koti zmanjšajo maksimalne pospeške, kar zmanjšuje obremenitve na astronavte in plovila.
    Pri velikih hitrostih in pravih trajektorijah smo opazili tudi "odbitek" od atmosfere zaradi vzgona, ki je deloval na plovilo.
    V modelu smo uporabili Newtonov zakon gravitacije, kvadratni zakon zračnega upora in vzgona, ter obstoječi atmosferski model NRLMSISE-00\cite{fig:vir5} za določanje gostote zraka.
    Simulacija ima omejitve in poenostavitve: dvodimenzionalni sistem, konstantna koeficienta upora in vzgona ter zanemarjanje termalnih učinkov in vpliva drugih nebesnih teles.
    V prihodnje bi lahko nadgradil model, da bi vseboval še temperaturo plovila, ter različne načine, kako se ščititi pred prevelikim segrevanjem.
    Lahko bi tudi iskal bolj optimalne poti skozi višjo atmosfero z aktivno kontrolo kapsule.
    Prostor bi lahko nadgradil z dodatno dimenzijo, lahko bi vpeljal bolj kompleksne vremenske pojave ipd.
    V prihodnosti nameravam nadgraditi model še z uporabo 3d koordinatnega sistema in Runge-Kutta integracijske metode.


    \pagebreak
    \section{Viri}
    \begin{thebibliography}{9}
        \bibitem[1]{fig:vir1}
        Wikipedia, Apollo command and service module. \\
        Dostopno na: \url{https://en.wikipedia.org/wiki/Apollo_command_and_service_module} (2025)
        \vspace{1cm}

        \bibitem[2]{fig:vir2}
        W. David Woods, Kenneth D. MacTaggart, and Frank O'Brien, Apollo 11 Day 9, part 2: Entry and Splashdown. \\
        Dostopno na: \url{https://www.nasa.gov/history/afj/ap11fj/27day9-entry.html} (2025)
        \vspace{1cm}

        \bibitem[3]{fig:vir3}
        Johnson, Breanna J., Mid-Lift-To-Drag Ratio Rigid Vehicle 6-DOF EDL Performance Using Tunable Apollo Powered Guidance. \\
        Dostopno na: \url{https://ntrs.nasa.gov/citations/20200001595} (2025)
        \vspace{1cm}

        \bibitem[4]{fig:vir4}
        Hillje, Ernest R., ``Entry Aerodynamics at Lunar Return Conditions Obtained from the Flight of Apollo 4 (AS-501),'' NASA TN D-5399, (1969). \\
        Dostopno na: \url{https://ntrs.nasa.gov/api/citations/19690029435/downloads/19690029435.pdf} (2025)
        \vspace{1cm}

        \bibitem[5]{fig:vir5}
        Wikipedia, NRLMSISE-00. \\
        Dostopno na: \url{https://en.wikipedia.org/wiki/NRLMSISE-00} (2025)

        \bibitem[6]{fig:vir6}
        Gary Johnson, Blunt Capsule Drag Data. \\
        Dostopno na: \url{https://exrocketman.blogspot.com/2012/08/blunt-capsule-drag-data.html} (2025)

    \end{thebibliography}

\end{document}
